\documentclass[11pt,a4paper]{article}

% 编码与字体支持
\usepackage[utf8]{inputenc}
\usepackage[T1]{fontenc}
\usepackage{amsmath, amssymb, amsthm}
\usepackage{geometry}
\usepackage{hyperref}

% 恢复正统的西里尔字母 Sha (Ш)，符合顶级数论期刊标准
\DeclareFontFamily{U}{wncyr}{}
\DeclareFontShape{U}{wncyr}{m}{n}{<->wncyr10}{}
\DeclareSymbolFont{mcy}{U}{wncyr}{m}{n}
\DeclareMathSymbol{\Sha}{\mathord}{mcy}{"58}

% 严格的数学算子声明
\DeclareMathOperator{\Gal}{Gal}
\DeclareMathOperator{\rank}{rank}

\geometry{margin=1in}

% 定理环境
\newtheorem{theorem}{Theorem}[section]
\newtheorem{lemma}[theorem]{Lemma}
\newtheorem{corollary}[theorem]{Corollary}
\newtheorem{remark}[theorem]{Remark}

\title{Arithmetic Geometry of Polynomial Curves: Decomposition of the Jacobian and the Tate-Shafarevich Group in Quadratic Twists of $y^2 = Q_7(x)$}
\author{
    Ruqing Chen \\
    GUT Geoservice Inc., Montr\'{e}al, QC, Canada \\
    \texttt{ruqing@hotmail.com}
}
\date{March 2026}

\newcommand{\githubrepo}{\url{https://github.com/Ruqing1963/Arithmetic-Geometry-Q7}}

\begin{document}

\maketitle

\begin{abstract}
The polynomials $Q_q(x) = x^q - (x-1)^q$ are of significant interest in computational number theory due to their prime-generating capacities under the Bateman-Horn heuristic \cite{BatemanHorn}. In this paper, we study the arithmetic geometry of the hyperelliptic curves $C_q: y^2 = Q_q(x)$ over $\mathbb{Q}$. For $q=7$, the curve $C_7$ has genus $g=2$. We prove that the splitting field of $Q_7(x)$ exhibits a strict $C_6$ Galois cyclic symmetry, allowing a non-trivial hyperelliptic involution. Through this involution, we explicitly construct a decomposition of the Jacobian \cite{CasselsFlynn}, obtaining a quotient elliptic curve $E_1$. Furthermore, by studying the family of square-free quadratic twists $E_1^{(D)}$, we perform a 2-descent analysis. At $D=167$, we demonstrate a failure of the Hasse principle by explicitly computing a non-trivial Tate-Shafarevich group $\Sha(E_1^{(167)}/\mathbb{Q})[2] \cong (\mathbb{Z}/2\mathbb{Z})^2$. By constructing the explicit quartic torsor and relying on the unconditional theorem of Kolyvagin, we provide a proof of this local-global obstruction that is independent of the Birch and Swinnerton-Dyer conjecture.
\end{abstract}

\smallskip
\noindent\textbf{Data and Code Availability.} All computational scripts (SageMath, Magma) and result data accompanying this paper are publicly available at \githubrepo.

\section{Introduction}

The polynomials $Q_q(x) = x^q - (x-1)^q$, which form the arithmetic core of the so-called ``Titan polynomials'', are of significant interest in computational number theory due to their prime-generating capacities under the Bateman-Horn heuristic \cite{BatemanHorn}. While our preceding numerical experiments focused on global prime enumeration, the Diophantine properties of the associated algebraic curves $C_q: y^2 = Q_q(x)$ remain relatively unexplored. While prime values require the polynomial evaluations to be prime integers, investigating the underlying algebraic geometry—specifically, the distribution of rational points and local-global principles—provides deeper insight into the arithmetic rigidity of these hyperelliptic structures.

When $q \ge 7$, the hyperelliptic curve $C_q$ reaches genus $g \ge 2$, making explicit computation of rational points and Selmer groups increasingly difficult due to the rapid growth of the conductor. In this paper, we study the arithmetic geometry of $C_7$ from a theoretical standpoint. We demonstrate that $Q_7(x)$ possesses a highly symmetric Galois structure, enabling algebraic dimension reduction (decomposition of the Jacobian). By systematically twisting the resultant genus-1 quotient curve, we identify an explicit arithmetic resonant point ($D=167$) where non-trivial elements of the Tate-Shafarevich group ($\Sha$) arise.

\section{Galois Symmetry and Decomposition of the Jacobian}

To study the Diophantine properties of the genus-2 curve $C_7: y^2 = x^7 - (x-1)^7$, we first analyze the Galois structure of the polynomial $Q_7(x)$.

\begin{lemma}[Cyclotomic Field and Galois Group]
The splitting field of $Q_7(x)$ over $\mathbb{Q}$ is isomorphic to the cyclotomic field $\mathbb{Q}(\zeta_7)$. Its Galois group $\Gal(Q_7/\mathbb{Q})$ is the cyclic group $C_6$.
\end{lemma}
\begin{proof}
Setting $Q_7(x) = 0$ over $\mathbb{C}$ is equivalent to $(\frac{x}{x-1})^7 = 1$. Excluding the trivial solution, the 6 roots are $x_k = (1 - \zeta_7^k)^{-1}$ for $k \in \{1, \dots, 6\}$, where $\zeta_7 = e^{2\pi i / 7}$. Thus, the splitting field is $\mathbb{Q}(\zeta_7)$. By classical Galois theory, its Galois group is isomorphic to $(\mathbb{Z}/7\mathbb{Z})^\times \cong C_6$. The existence of this abelian structure ($|C_6| = 2 \times 3$) guarantees a normal subgroup of index 2, allowing for a decomposition of the corresponding geometric space.
\end{proof}

\begin{theorem}[Decomposition of the Jacobian via Hyperelliptic Involution]
The hyperelliptic curve $C_7: y^2 = Q_7(x)$ of genus $g=2$ admits a non-trivial involution $\sigma$, yielding a quotient elliptic curve $E_1$ of genus $g=1$.
\end{theorem}
\begin{proof}
Expanding $Q_7(x)$, we observe an absolute reflection symmetry:
$$Q_7(1-x) = (1-x)^7 - (-x)^7 = (1-x)^7 + x^7 = Q_7(x)$$
This induces an involution $\sigma: (x,y) \mapsto (1-x, y)$ on $C_7$. We define the invariant $t = x(1-x) = x - x^2$. Rewriting $Q_7(x)$ in terms of $t$:
$$Q_7(x) = 1 - 7(x-x^2) + 14(x-x^2)^2 - 7(x-x^2)^3$$
Substituting $t$, we obtain a cubic polynomial:
$$Q_7(t) = -7t^3 + 14t^2 - 7t + 1$$
The quotient curve under $\sigma$ is $y^2 = -7t^3 + 14t^2 - 7t + 1$. By applying the affine transformation $X = -7t$ and $Y = 7y$, we derive the Weierstrass standard form of the quotient curve $E_1$:
$$E_1: Y^2 = X^3 + 14X^2 + 49X + 49$$
This explicit construction demonstrates the decomposition of the Jacobian: $Jac(C_7) \sim E_1 \times A$, where $A$ is the corresponding Prym variety \cite{CasselsFlynn}.
\end{proof}

\section{Quadratic Twists and the 2-Selmer Group}

To investigate the local-global obstructions associated with $E_1$, we introduce square-free quadratic twists $D \in \mathbb{Z}$, constructing the family:
$$E_1^{(D)}: Y^2 = X^3 + 14DX^2 + 49D^2X + 49D^3$$
Twisting by $D$ modifies the conductor and the local reduction properties at primes dividing $D$, inducing variations in the 2-Selmer group $Sel_2(E_1^{(D)}/\mathbb{Q})$. For lower values of $D \le 47$, computational 2-descent shows that the 2-Selmer rank strictly aligns with the analytic rank, confirming the triviality of $\Sha(E_1^{(D)}/\mathbb{Q})[2]$. 

\begin{remark}
In computational arithmetic geometry, heuristic bounds for $Sel_2$ or analytic rank can sometimes violate the Parity Conjecture if the exact generators or the true rank are not rigorously verified. Therefore, identifying a true failure of the Hasse principle requires rigorous verification of both the analytic rank and the explicit 2-covering spaces.
\end{remark}

\section{Failure of the Hasse Principle at \texorpdfstring{$D=167$}{D=167}}

Through systematic heuristic scanning via analytic $L$-functions (which conditionally rely on the Birch and Swinnerton-Dyer conjecture to estimate the order of $\Sha$), we isolated a highly probable instance of local-global failure at $D=167$. Let $E = E_1^{(167)}$ for brevity. Substituting $D=167$ into the twist family gives the explicitly computed curve:
$$E: Y^2 = X^3 + 2338X^2 + 1366561X + 228215687$$

To bypass the reliance on the BSD conjecture, we provide a proof constructed via explicit 2-descent and Kolyvagin's unconditional theorem.

\begin{theorem}
The twisted elliptic curve $E$ has a non-trivial Tate-Shafarevich group. Specifically, its 2-torsion subgroup is isomorphic to the Klein four-group:
$$\Sha(E/\mathbb{Q})[2] \cong (\mathbb{Z}/2\mathbb{Z})^2$$
\end{theorem}
\begin{proof}
We evaluate the analytic rank via the $L$-series. Computations via modular symbols confirm that the central value $L(E, 1) \neq 0$. By the unconditional theorem of Kolyvagin \cite{Kolyvagin} and Gross-Zagier \cite{GrossZagier}, a non-zero central $L$-value implies that the Mordell-Weil group over $\mathbb{Q}$ has rank exactly $0$. Furthermore, testing for rational roots reveals that the polynomial $X^3 + 2338X^2 + 1366561X + 228215687$ has no roots in $\mathbb{Q}$, meaning the rational 2-torsion subgroup is trivial: $E(\mathbb{Q})[2] = 0$.

Concurrently, we compute the 2-Selmer group $Sel_2(E/\mathbb{Q})$ via an explicit 2-descent algorithm. Testing the 2-coverings for local solubility over $\mathbb{R}$ and all $p$-adic fields $\mathbb{Q}_p$ yields an $\mathbb{F}_2$-dimension:
$$\dim_{\mathbb{F}_2} Sel_2(E/\mathbb{Q}) = 2$$

Applying the fundamental short exact sequence for the descent:
$$0 \to E(\mathbb{Q})/2E(\mathbb{Q}) \to Sel_2(E/\mathbb{Q}) \to \Sha(E/\mathbb{Q})[2] \to 0$$
Given that $\rank(E(\mathbb{Q})) = 0$, the Mordell-Weil group consists entirely of its torsion subgroup $E(\mathbb{Q})_{tors}$. Since there are no rational 2-torsion points ($E(\mathbb{Q})[2] = 0$), the order of this finite abelian group must be odd. For a finite abelian group of odd order, the multiplication-by-2 map is an automorphism, implying $2E(\mathbb{Q}) = E(\mathbb{Q})$. Therefore, the quotient group $E(\mathbb{Q})/2E(\mathbb{Q})$ is trivial ($0$-dimensional). The exactness of the sequence thus forces an isomorphism:
$$Sel_2(E/\mathbb{Q}) \cong \Sha(E/\mathbb{Q})[2]$$

Since $\dim_{\mathbb{F}_2} Sel_2 = 2$, the Selmer group contains exactly $2^2 = 4$ elements. One of these elements corresponds to the trivial principal homogeneous space (isomorphic to $E$ itself), which possesses a global rational point (the point at infinity $\mathcal{O}$). The remaining three elements are non-trivial principal homogeneous spaces.

For instance, computing the explicit 2-descent via classical invariant theory and Cassels' algorithm \cite{Cassels} (implemented in the \texttt{Magma} computational algebra system \cite{Magma}) yields the following explicitly locally soluble quartic torsor:
$$v^2 = 12u^4 - 163u^3 - 83u^2 + 2120u + 2243$$
This explicit principal homogeneous space is locally soluble over $\mathbb{R}$ and all completions $\mathbb{Q}_p$, but possesses absolutely no global rational points over $\mathbb{Q}$. Its existence provides a constructive proof for a strict failure of the Hasse principle, independent of the BSD conjecture.
\end{proof}

\section{Conclusion}

By shifting the analysis from global prime enumeration to the arithmetic geometry of hyperelliptic curves, we have demonstrated the decomposability of the Jacobian of $y^2 = x^7 - (x-1)^7$ via its inherent Galois symmetries. Through quadratic twists, we rigorously established the existence of non-trivial elements in the Tate-Shafarevich group at $D=167$ via explicit 2-descent and Kolyvagin's theorem. 

Furthermore, comparing these results with the $q=5$ case offers broader arithmetic insights. The elliptic curves associated with $Q_5(x)$ typically exhibit positive Mordell-Weil rank at very small quadratic twists (e.g., $D=1, 2, 3$), indicating an abundance of rational points. In stark contrast, the dimension-reduced elliptic curve for $Q_7(x)$ exhibits profound arithmetic rigidity, maintaining rank zero and structural smoothness until the singular resonance at $D=167$. This suggests that as the degree of the underlying Titan polynomial increases, its arithmetic geometric structures become increasingly resistant to producing rational points, culminating in cohomological obstructions. Consequently, while the Bateman-Horn heuristic predicts prime densities globally, the extreme algebraic rigidity and local-global failures embedded within these hyperelliptic curves impose profound geometric constraints that may inherently suppress the arithmetic density of prime constellations generated by higher-degree polynomials.

% ================= 参考文献 (Bibliography) =================
\begin{thebibliography}{9}

\bibitem{BatemanHorn}
P. T. Bateman and R. A. Horn,
\textit{A heuristic asymptotic formula concerning the distribution of prime numbers},
Mathematics of Computation, 16(79):363--367, 1962.

\bibitem{Cassels}
J. W. S. Cassels,
\textit{Lectures on Elliptic Curves},
London Mathematical Society Student Texts 24, Cambridge University Press, 1991.

\bibitem{CasselsFlynn}
J. W. S. Cassels and E. V. Flynn,
\textit{Prolegomena to a Middlebrow Arithmetic of Curves of Genus 2},
London Mathematical Society Lecture Note Series 230, Cambridge University Press, 1996.

\bibitem{GrossZagier}
B. H. Gross and D. B. Zagier,
\textit{Heegner points and derivatives of L-series},
Inventiones mathematicae, 84(2):225--320, 1986.

\bibitem{Kolyvagin}
V. A. Kolyvagin,
\textit{Finiteness of $E(\mathbb{Q})$ and $\Sha(E,\mathbb{Q})$ for a subclass of Weil curves},
Izvestiya Akademii Nauk SSSR. Seriya Matematicheskaya, 54(3):522--559, 1990.

\bibitem{Magma}
W. Bosma, J. Cannon, and C. Playoust,
\textit{The Magma algebra system. I. The user language},
Journal of Symbolic Computation, 24(3-4):235--265, 1997.

\end{thebibliography}

\end{document}